\section{问题重述}

\subsection{背景}

在生鲜商超中，蔬菜类商品的保鲜期较短，品相随销售时间增加而变差，大部分品种当日未售出隔日就无法再售。商超通常根据各商品的历史销售和需求情况每天进行补货，进货交易时间在凌晨 $3$ 时 $00$ 分至 $4$ 时 $00$ 分之间，商家在不确切知道具体单品和进货价格的情况下需做出当日补货决策。蔬菜的定价一般采用成本加成定价方法，对运损和品相变差的商品通常进行打折销售。

\subsection{问题描述}

附件 1 给出了某商超经销的 6 个蔬菜品类的商品信息（共 251 个单品），附件 2 和附件 3 分别给出了 2020 年 7 月 1 日至 2023 年 6 月 30 日的销售流水明细与批发价格数据，附件 4 给出了各商品的近期损耗率。根据附件和实际情况需要解决以下四个问题，

\textbf{问题一}要求分析蔬菜各品类及单品销售量的分布规律及相互关系。

\textbf{问题二}要求以品类为单位做补货计划，分析各蔬菜品类的销售总量与成本加成定价的关系，并给出各蔬菜品类未来一周（2023 年 7 月 1 日至 7 日）的日补货总量和定价策略，使得商超收益最大。

\textbf{问题三}要求制定单品补货计划，可售单品总数控制在 $27$ 至 $33$ 个，各单品订购量满足最小陈列量 $2.5$ 千克的要求。根据 2023 年 6 月 24 日至 30 日的可售品种，给出 7 月 1 日的单品补货量和定价策略，在尽量满足市场对各品类蔬菜商品需求的前提下，使得商超收益最大。

\textbf{问题四}要求提出商超为更好制定补货和定价决策还需要采集的相关数据，并说明数据对解决问题的帮助、意见和理由。

\section{问题分析}

\subsection{问题一分析}

\textbf{首先}，问题一属于描述推断型任务，是后续建模的认知基础。\textbf{其次}，需要从两个层面展开分析，品类层面关注销售分布特征和品类间关联结构，单品层面关注销售分散度和单品聚类。\textbf{然后}，一个被 5 篇已有文献忽略的关键问题是，共同时间趋势（季节性、节假日效应）会系统性夸大品类间的 Spearman 相关系数，因此在计算关联前必须先对销量序列做时间序列分解并剥离趋势。\textbf{最后}，相关不等同于因果，需要引入 Granger 因果检验区分因果方向。

\subsection{问题二分析}

\textbf{首先}，问题二属于优化决策型任务，核心是收益最大化。\textbf{其次}，需要处理三个不确定性来源，需求预测的不确定性（报童模型）、批发价格的不确定性（预测区间而非点估计）、价格弹性估计的不确定性。\textbf{然后}，成本加成定价中的加成率不应作为固定外生参数，需要内生于优化模型。\textbf{最后}，6 个品类的需求并非完全独立（问题一的结论），但品类间的残差相关强度中等（平均 $|\rho|=0.266$），在优化模型中可近似为独立决策，但需要在鲁棒性分析中考虑相关性的影响。

\subsection{问题三分析}

\textbf{首先}，问题三在问题二的基础上从品类级下沉到单品级。\textbf{其次}，核心约束有三个，单品总数限制（$27$ 至 $33$ 个）、最小陈列量（$2.5$ kg）、品类覆盖要求。\textbf{然后}，选品、补货量、定价三者是耦合的，定价影响需求，需求决定补货量，补货量约束选品范围，因此需要联合优化而非分步求解。\textbf{最后}，品类最低覆盖数应根据各品类在前期（6 月 24 日至 30 日）的实际可售品种数等比例分配。

\subsection{问题四分析}

问题四是开放性建议题，需要基于问题一至问题三的实际建模过程中暴露的问题，提出具体可行的数据采集方案。每条建议需对应前三问中发现的具体建模缺陷，而非泛泛而谈。

\section{数据预处理}

\subsection{数据来源与结构}

附件 1 包含 251 个蔬菜单品信息（单品编码、名称、品类编码、品类名称），分为 6 个品类，花叶类（100 个单品）、食用菌（72）、花菜类（45）、水生根茎类（19）、辣椒类（10）、茄类（5）。附件 2 包含 $878,503$ 条销售流水记录，时间跨度为 2020 年 7 月 1 日至 2023 年 6 月 30 日（$1,085$ 天），字段包括销售日期、扫码时间、单品编码、销量（kg）、销售单价（元/kg）、销售类型（销售/退货）、是否折扣销售。附件 3 包含 $55,982$ 条批发价格记录，时间跨度与附件 2 一致。附件 4 包含品类级和单品级近期损耗率。

\subsection{数据清洗}

四个附件均不存在缺失值。在此基础上做三项清洗。

\textbf{退货处理}。附件 2 含 $461$ 条退货记录（占总量 $0.05\%$），销量字段为负值。退货原因包括品相问题和顾客退货等，与正常需求模式无关，直接剔除。保留 $878,042$ 条正常销售记录进入后续分析。

\textbf{品类聚合}。将单品级销售流水按日期和品类汇总，得到 $6$ 品类 $\times$ $1,085$ 天的日销量矩阵。批发价按同样方式聚合，对没有批发价记录的日期使用前向填充，蔬菜批发价通常连续定价，日间不会剧烈跳变。

\textbf{单品持续性筛选}。$246$ 个有销售记录的单品中，$87$ 个（$35.4\%$）的有效销售天数不足 $30$ 天。这些低持续性单品在时间序列层面的信号太弱，不是需求为零就是偶发性上架，直接纳入聚类会引入噪声。单品层面的分析限定在销售持续性好的 $116$ 个单品（$\geq 90$ 天销售记录）。

\subsection{基本统计特征}

品类销量分布极不均衡。花叶类三年累计销量 $198,521$ kg（$\mathbf{42.2\%}$），辣椒类 $91,589$ kg（$\mathbf{19.4\%}$），食用菌 $76,087$ kg（$16.2\%$），花菜类 $41,766$ kg（$8.9\%$），水生根茎类 $40,581$ kg（$8.6\%$），茄类仅 $22,432$ kg（$\mathbf{4.8\%}$）。品类不平衡比为 $8.8$（花叶类/茄类）。

全品类日总销量均值 $434$ kg，标准差 $208$ kg（日变异系数约 $0.48$），最大单日销量 $2,484$ kg。8 月是全年销售旺季（月均 $51,859$ kg），5 月至 6 月为低谷（月均约 $30,000$ kg），季节性振幅约 $1.7$ 倍。

单品销量高度集中。246 个有销售的单品按三年总销量降序排列，Lorenz 曲线显示 Gini 系数 $\mathbf{0.793}$，前 $20\%$ 单品（约 $50$ 个）贡献 $83.9\%$ 总销量。这条集中度曲线对后续定价和选品有直接含义，少量明星单品驱动绝大部分收入，长尾单品仅因品类多样性才需要纳入考虑。
